\scanpage{42}
% archive sheet number: 41

\[
 X_{\mathrm{cpct}}\simeq\widetilde R_{\mathrm{top\,cpct}}
\]

cpct est la topologie sur $R$ pour laquelle $X_i\longrightarrow X$ est couv.
ssi pour tt anneau \inserted{local} artinien \deleted{\unsure{local}}
(ou \unsure{artinien|artin.|anneau} de dim 0) \deleted{\unsure{absolument quelconque}} $O$, $X_i(O)\to X(O)$ est épimorphique

encore un topos non cohérent (bien que sous-topos d'un topos cohérent)

\textcircled{\scriptsize\uline{11}} \deleted{\unsure{Pseudo}}pseudo-corps $=$ \deleted{\unsure{compact réduit}}
\inserted{Anneau \inserted{réduit} (local de dim 0)}

\[
 A^*\amalg_{\widehat R} A\xrightarrow{\mathrm{incl},\,0_A}A\quad\text{est épi}
\]

$X_{\mathrm{ps.corps}}\simeq\widetilde R_{\deleted{\unsure{\mathrm{top}}}\mathrm{cons}}$\quad cons, top. sur $R$
où les familles couvrantes sont celles qui ont \deleted{\unsure{une section sur toute fibre}} des sections sur chaque fibre.

$\simeq\widetilde R_{\rm cons}$\quad $R_{\rm cons}$ la catégorie des schémas compacts réduits associés aux schémas de t.f. sur $\mathbb Z$, avec top. de Zariski.

$\widetilde R_{\rm cons}$ formé des foncteurs contravariants qui commutent aux \unsure{sommes|lim} finies

\deleted{\textcircled{\scriptsize 12} corps $=$ compact \deleted{\unsure{intègre}}\inserted{intègre} $=$ ps-corps intègre
$=$ pseudo-corps local \deleted{\unsure{(ou quasi-intègre) ou bien local et réduit ou quasi-intègre}}}

% slanted
\margin{clair}\quad\margin{à prouver}

\deleted{\uline{NB} les ss-topos définis dans 1) à 11) \deleted{\unsure{sont}} ont assez de pts, donc sont définis par leurs points.
Il n'est pas clair qu'il en soit nécessairement ainsi pour les intersections de plusieurs de ces sous topos.
Donc pour vérifier que deux intersections sont égales, \deleted{\unsure{il}}\ ou ont des rel d'inclusion déterminées,
il n'est pas sûr qu'il suffira de le vérifier pour leurs pts, i.e. pour les structures ensemblistes.
-- Sauf si on vérifie à chaque fois que les int. envisagées ont assez de pts !}

\deleted{Vérifier donc que un Anneau \unsure{compact|cpct} \deleted{intègre} A est \deleted{\unsure{intègre}} \inserted{local hens.}
ps corps local A est \unsure{intègre} \inserted{corps alg.} ou cpct intègre,}

\deleted{\unsure{D)|2)} La définition signifie que $A:R\to\mathcal E$ transforme sommes \deleted{familles} couvrantes pour ``top. cpct'' en
\deleted{\unsure{fibrés à sommes finies}}, et qu'il transforme familles couvrantes pour ``top. cpct'' en familles couvrantes,
Montrons qu'une famille $X_i\to X$ qui \deleted{\unsure{de}} est couvrante pour topo discrète au dessus de tt $x\in X$
(i.e. $X_i(k)\to X(k)$ est épi pour tt corps $k$) \deleted{\unsure{est transformée}} en famille épi.
Réc. sur $\dim X$ assez évidente.}
